\centerline{\bf \Huge Начало комбинаторики II}

\begin{flushright}
        \scriptsize{Никогда не сдаваться... Встать, когда все рухнуло — вот настоящая сила. Наруто}
\end{flushright}

\problem{\textbf{1}} Начальник хочет побывать в Мадриде, Сеуле, Токио и Минске, но ещё не решил, в какой
последовательности. Сколько перед ним различных вариантов выбора маршрута?

\problem{\textbf{2}} Сколько двузначных чисел можно составить из цифр 0, 1, 2, 3, 4, 5, 6?

\problem{\textbf{3}} Сколько существует трёхзначных чисел, сумма цифр которых не превосходит 4?

\problem{\textbf{4}} Сколько трёхзначных чисел можно составить из цифр 0, 1, 2, 3, 4, 5, 6?

\problem{\textbf{5}}  Есть две белые, две красные и две розовые гвоздики. Сколькими способами их можно
расставить в три вазы так, чтобы в каждой вазе стояли по две гвоздики разного цвета?

\problem{\textbf{6}}  Сколькими способами из цифр 1, 2, 3, 4 можно составить
число, кратное 6? При составлении числа каждую цифру можно использовать один раз или не
использовать совсем.

\problem{\textbf{7}} Каких пятизначных чисел больше — тех, в записи которых есть цифра 7, или тех, в записи которых ее нет? И на сколько?